\documentclass{zhvt-classic}

% 使用本地源流明体（GenRyuMin2 TW）按文件路径加载，不依赖系统字体名
% 版框（内框／界栏／书口侧栏／鱼尾）改用与版心印章同一印泥棕色：
%   印章图片 yshs.png 实色 = RGB(135,68,52) = HTML 874434
\usepackage{xcolor}
\usepackage{tikz}
% sealbrown 双注册：xcolor 供 TikZ/pgf，L3 供 cls 版框
\definecolor{sealbrown}{HTML}{874434}
\ExplSyntaxOn
\color_set:nnn { sealbrown } { HTML } { 874434 }
\ExplSyntaxOff

\zhvtset{
  grid color       = sealbrown,
  main font        = GenRyuMin2TW-SB,
  main font path   = C:/Users/Administrator/texmf/fonts/opentype/public/GenRyuMin2TW/,
jiazhu font      = GenRyuMin2TW-M,
  jiazhu font path = C:/Users/Administrator/texmf/fonts/opentype/public/GenRyuMin2TW/,
  % 页码框容纳 4 字：正常页「九十七」+前置空 1 格=4 行恰好；
  % 100+ 页「一百零一」共 5 行超出，末字「一/二/三」被框裁掉（2026-10-04 实测），故放宽到 5
  shukou page number chars = 5,
  % 书口书名框高从默认 6 字减到 4 字，把省出的 2 字高拨给 chapter 框，
  % 使 6 字书名+卷名放进 9 字高的 chapter 框后不再压线（2026-10-05）
  shukou title chars = 4,
}

\title{史记}[魏其武安侯列传]
\author{汉·司马迁}
\maker{闲叟漫记·重製}


% ===== 索隐述赞样式（2026-11-08 四次定稿 · 愚大人亲手改定参数）=====
% ⚠ 竖排前提：自定义字体家族必须带 RawFeature={vertical:+vert}；
%   simhei.ttf 加 +vhal 会整体加载失败静默回退，只写 +vert；simfang 同理。
\newCJKfontfamily{\shuzanfont}[Path=C:/Windows/Fonts/,RawFeature={vertical:+vert}]{simhei.ttf}
\newCJKfontfamily{\shuzanbody}[Path=C:/Windows/Fonts/,RawFeature={vertical:+vert}]{simfang.ttf}
% 题头：152x30pt 圆角 6pt 棕底白字；baseline (0,2) 使整框落入界行格子不压版框线；
% 白字节点 x 不补偿、y +8.2pt（像素级墨迹实测 dx/dy≈0，L/R 留白约 20pt）。
\newcommand{\shuzanhead}[1]{%
  \begin{tikzpicture}[baseline={(0,2)}]
    \fill[sealbrown, rounded corners=6pt]
      (-76pt,-15pt) rectangle (76pt,15pt);
    \node[text=white, inner sep=0pt, anchor=center] at (0pt,8.2pt)
      {\shuzanfont\fontsize{23}{23}\selectfont #1};
  \end{tikzpicture}%
}
% ★ 修正书口卷名位置：chapter 框起始位置改回 cls 默认值，避免全局影响。
%   只在每个卷首页通过 \juan 命令里的空命令微调该页位置（安全，不会跨页污染）。
\ExplSyntaxOn
\AtBeginDocument{
  \dim_set:Nn \l__zhvt_shukou_chapter_title_start_dim
    { \l__zhvt_shukou_title_height_dim + \l__zhvt_font_size_dim * 2 }
  \dim_set:Nn \l__zhvt_shukou_page_number_start_dim
    { \l__zhvt_shukou_chapter_title_start_dim + \l__zhvt_shukou_chapter_title_height_dim }
}
\ExplSyntaxOff

\begin{document}

% ★★ 版心（书口）排布覆写：把书名放到鱼尾之下（2026-10-04）
%
%   cls 默认版心自上而下：书名（title 静态框，在鱼尾之上）→ 鱼尾 → 篇名（chapter 动态框）→ 页码
%   本模板改为：          鱼尾（位置不动）→ 书名 → 篇名 → 页码分割线 → 页码
%
%   手法：把 cls 放在 title 静态框里的书名，改放进 chapter 动态框（鱼尾之下）；
%         title 静态框留空。书名取 \l__zhvt_title_str，即 \title{...} 的第一参数，
%         所以换书时只改 \title，不必在此重复写书名。
%
%   ⚠ 三条必须遵守的写法（2026-10-04 实测）：
%     ① 汉字必须用 \hbox{...} 包住。flowfram 会把框内容装进一个宽 14.4mm（= 一个 line_sep
%        = 书口栏宽）的纵向 minipage 并按该宽度折行，而整页靠 PDF /Rotate 90 显示、
%        框内文字本身不变向 —— 不包 \hbox 时「围炉夜话」4 字会被逐字折成 4 行，
%        在版面上就成了横排 4 字并横向溢出版框。
%     ② 竖排下**一个 \hbox 就是一整列**。书名与篇名要同处一列，必须放进同一个 \hbox；
%        分成两个 \hbox 时，后一个会被当成新的一列、甩到隔壁栏去
%        （实测「序」跑到书口栏外侧 14.4mm 的正文列上）。中间的 \quad 在框内是横向空白，
%        显示出来正好是书名与篇名之间的纵向间隔。
%     ③ ★ 内容盒要写成 **与框等宽的定宽盒** `\hbox to \l__zhvt_shukou_chapter_title_height_dim{…\hfil}`。
%        实测：chapter 框是把内容盒沿纵向**居中**摆放的，内容宽度一变，书名就跟着漂：
%        自然宽度的 \hbox（书名 38.4mm）落在 x=157.2mm，加篇名后（57.6mm）落在 x=152.4mm，
%        目录页（无篇名）与正文页的书名就差半格。做成等宽盒后，无论有没有篇名，
%        内容固定从框顶起排，实测两页均落在 x=142.8mm，不再漂。
\ExplSyntaxOn
% 把书口 chapter 框（书名+篇名）起始位置下移 1 字高，避免 4 字卷名顶到鱼尾
% ⚠ cls 的 \__zhvt_calculate_dimensions: 在 \AtBeginDocument 里执行，会重算该 dim；
%   必须把覆写挂到 begin document 之后，否则会被二次覆盖（2026-10-05 实测）。
% (旧方案已移到导言区，直接修改 chapter 框起始位置)
% 卷标题命令 #1=书口篇名(鱼尾) #2=正文卷名（含卷序，卷名内空格用 ~）
%
% ★★ 目录页两项修订（2026-10-05，愚大人指出）★★
%   ① 页码起点统一对齐：
%      原状 —— 目录项写成 \par\hspace{4em}卷名\hspace{4em}页码，页码起点 = 卷名实宽 + 4em。
%      卷名实宽随字数变：「卷一~梁惠王上」6 字 +~、「卷七~离娄上」5 字 +~，
%      差 1 个字位 → 卷七~卷十 4 条的页码整体比其余 10 条高 27.21pt（一个字），看着上下不齐。
%      改法 —— 卷名补全角空格（宽 = 1 字 = 27.215pt）补足到 **7 字符位**再写进目录，
%      14 条页码起点遂全部落在同一条竖线上（起点对齐；末字随页码字数自然参差）。
%      ⚠ 只对目录里的副本 \l_tmpc_tl 补，正文卷名仍用 \l_tmpa_tl 原样 —— 正文像素不受影响。
%   ② 卷名整体上移半字位：前导 \hspace 由 4em 改 3.5em。
%      实测原状目录块在版框内偏下（天头留白 139.6pt / 地脚留白 110.7pt），
%      上移 0.5em = 13.6pt 后上下基本均分，视觉居中。
\NewDocumentCommand { \juan } { mm } {
	\clearpage
	\tl_set:Nn \l__zhvt_chapter_title_tl { #1 }
	\tl_set:Nn \l_tmpa_tl { #2 }
	\tl_set_eq:NN \l_tmpc_tl \l_tmpa_tl
	\int_step_inline:nn { 7 - \str_count:n { #2 } }
	  { \tl_put_right:Nn \l_tmpc_tl { 　 } }
	\tl_set:Nx \l_tmpb_tl { \zhnum { page } }
	\tl_put_left:Nn \l_tmpb_tl { \hspace { 4em } }
	\tl_put_left:Nn \l_tmpb_tl { \l_tmpc_tl }
	\tl_put_left:Nn \l_tmpb_tl { \par \hspace { 3.5em } }
	\addtocontents { toc } { \l_tmpb_tl }
	\skip_horizontal:n { \l__zhvt_font_size_dim * 4 }
	\l_tmpa_tl
	\par
	\thispagestyle { empty }
}
\cs_set_protected:Npn \__zhvt_set_shukou_content: {
	% ★ 鱼尾上方留白（title 静态框）嵌入闲叟钤印：
	%   整页 /Rotate 90，故此处把图用 \rotatebox{90} 转正为竖章；
	%   尺寸卡在版心条幅内（书内约 7mm 宽 × 33.8mm 高，留余量不压鱼尾）。
	\setstaticcontents * { title } {
		\vbox_to_ht:nn { \l__zhvt_baseline_skip_dim } {
			\vskip 14.2pt
			\hbox_to_wd:nn { \l__zhvt_shukou_title_height_dim } {
				\hfil
				\rotatebox [ origin = c ] { 90 } { \includegraphics [ width = 7mm , height = 33.8mm ] { yshs.png } }
				\hfil
			}
		}
	}
% ★ 版心（书口）统一改法（2026-11-16 愚大人定）：18pt＋框内居中＋书口文字同版框色
%   ① 原 27.2pt 时 7 字长篇名（如「廉颇蔺相如列传」）末字溢出 chapter 框、压住页码；
%      改 \fontsize{18pt}{18pt}，前后 \hfil 令字串在框内居中（不随篇名字数漂移）。
%   ② \vbox{\vskip 8.2pt} 定量下移，抵消缩字号后栏宽方向 8.2pt 上偏（像素实测 y 中心 543.3 = 原版）。
%      ⚠ \raisebox / \vspace* 在 flowfram 框内无效；换字号须同步换补偿：16pt→10.1、19pt→7.4。
%   ③ chapter 与 page~number 均置 \color{sealbrown}（HTML 874434），与版框／界行同色；正文仍黑。
%   ⚠ 页码字位、钤印、鱼尾、界行、正文一律未动；五部合订本同此口径。
	\setdynamiccontents * { chapter }
	{ \color{sealbrown} \vbox { \vskip 8.2pt \hbox to \l__zhvt_shukou_chapter_title_height_dim
		{ \hfil \fontsize{18pt}{18pt}\selectfont \l__zhvt_title_str \quad \l__zhvt_chapter_title_tl \hfil } } }
	\setdynamiccontents * { page~number } { \color{sealbrown} \quad \zhnum { page } }
}
\__zhvt_set_shukou_content:
\ExplSyntaxOff

\maketitle

% 单篇，无多卷目录，不渲染目录页

\juan{魏其武安侯列传}{卷一百零七~魏其武安侯列传第四十七}
\hspace*{2em}魏其侯窦婴者。孝文后从兄子也。父世观津人。喜宾客。孝文时。婴为吴相。病免。孝景初即位。为詹事。\\
\hspace*{2em}梁孝王者。孝景弟也。其母窦太后爱之。梁孝王朝。因昆弟燕饮。是时上未立太子。酒酣。从容言曰。﹁千秋之后传梁王。﹂太后欢。窦婴引卮酒进上。曰。﹁天下者。高祖天下。父子相传。此汉之约也。上何以得擅传梁王。﹂太后由此憎窦婴。窦婴亦薄其官。因病免。太后除窦婴门籍。不得入朝请。\\
\hspace*{2em}孝景三年。吴楚反。上察宗室诸窦毋如窦婴贤。乃召婴。婴入见。固辞谢病不足任。太后亦惭。於是上曰。﹁天下方有急。王孙宁可以让邪。﹂乃拜婴为大将军。赐金千斤。婴乃言袁盎。栾布诸名将贤士在家者进之。所赐金。陈之廊庑下。军吏过。辄令财取为用。金无入家者。窦婴守荥阳。监齐赵兵。七国兵已尽破。封婴为魏其侯。诸游士宾客争归魏其侯。孝景时每朝议大事。条侯。魏其侯。诸列侯莫敢与亢礼。\\
\hspace*{2em}孝景四年。立栗太子。使魏其侯为太子傅。孝景七年。栗太子废。魏其数争不能得。魏其谢病。屏居蓝田南山之下数月。诸宾客辩士说之。莫能来。梁人高遂乃说魏其曰。﹁能富贵将军者。上也。能亲将军者。太后也。今将军傅太子。太子废而不能争。争不能得。又弗能死。自引谢病。拥赵女。屏间处而不朝。相提而论。是自明扬主上之过。有如两宫螫将军。则妻子毋类矣。﹂魏其侯然之。乃遂起。朝请如故。\\
\hspace*{2em}桃侯免相。窦太后数言魏其侯。孝景帝曰。﹁太后岂以为臣有爱。不相魏其。魏其者。沾沾自喜耳。多易。难以为相。持重。﹂遂不用。用建陵侯卫绾为丞相。\\
\hspace*{2em}武安侯田蚡者。孝景后同母弟也。生长陵。魏其已为大将军后。方盛。蚡为诸郎。未贵。往来侍酒魏其。跪起如子姓。及孝景晚节。蚡益贵幸。为太中大夫。蚡辩有口。学槃盂诸书。王太后贤之。孝景崩。即日太子立。称制。所镇抚多有田蚡宾客计筴。蚡弟田胜。皆以太后弟。孝景后三年封蚡为武安侯。胜为周阳侯。\\
\hspace*{2em}武安侯新欲用事为相。卑下宾客。进名士家居者贵之。欲以倾魏其诸将相。建元元年。丞相绾病免。上议置丞相。太尉。籍福说武安侯曰。﹁魏其贵久矣。天下士素归之。今将军初兴。未如魏其。即上以将军为丞相。必让魏其。魏其为丞相。将军必为太尉。太尉。丞相尊等耳。又有让贤名。﹂武安侯乃微言太后风上。於是乃以魏其侯为丞相。武安侯为太尉。籍福贺魏其侯。因吊曰。﹁君侯资性喜善疾恶。方今善人誉君侯。故至丞相。然君侯且疾恶。恶人众。亦且毁君侯。君侯能兼容。则幸久。不能。今以毁去矣。﹂魏其不听。\\
\hspace*{2em}魏其。武安俱好儒术。推毂赵绾为御史大夫。王臧为郎中令。迎鲁申公。欲设明堂。令列侯就国。除关。以礼为服制。以兴太平。举適诸窦宗室毋节行者。除其属籍。时诸外家为列侯。列侯多尚公主。皆不欲就国。以故毁日至窦太后。太后好黄老之言。而魏其。武安。赵绾。王臧等务隆推儒术。贬道家言。是以窦太后滋不说魏其等。及建元二年。御史大夫赵绾请无奏事东宫。窦太后大怒。乃罢逐赵绾。王臧等。而免丞相。太尉。以柏至侯许昌为丞相。武彊侯庄青翟为御史大夫。魏其。武安由此以侯家居。\\
\hspace*{2em}武安侯虽不任职。以王太后故。亲幸。数言事多效。天下吏士趋势利者。皆去魏其归武安。武安日益横。建元六年。窦太后崩。丞相昌。御史大夫青翟坐丧事不办。免。以武安侯蚡为丞相。以大司农韩安国为御史大夫。天下士郡诸侯愈益附武安。\\
\hspace*{2em}武安者。貌侵。生贵甚。又以为诸侯王多长。上初即位。富於春秋。蚡以肺腑为京师相。非痛折节以礼诎之。天下不肃。当是时。丞相入奏事。坐语移日。所言皆听。荐人或起家至二千石。权移主上。上乃曰。﹁君除吏已尽未。吾亦欲除吏。﹂尝请考工地益宅。上怒曰。﹁君何不遂取武库。﹂是后乃退。尝召客饮。坐其兄盖侯南乡。自坐东乡。以为汉相尊。不可以兄故私桡。武安由此滋骄。治宅甲诸第。田园极膏腴。而市买郡县器物相属於道。前堂罗锺鼓。立曲旃。后房妇女以百数。诸侯奉金玉狗马玩好。不可胜数。\\
\hspace*{2em}魏其失窦太后。益疏不用。无势。诸客稍稍自引而怠傲。唯灌将军独不失故。魏其日默默不得志。而独厚遇灌将军。\\
\hspace*{2em}灌将军夫者。颍阴人也。夫父张孟。尝为颍阴侯婴舍人。得幸。因进之至二千石。故蒙灌氏姓为灌孟。吴楚反时。颍阴侯灌何为将军。属太尉。请灌孟为校尉。夫以千人与父俱。灌孟年老。颍阴侯彊请之。郁郁不得意。故战常陷坚。遂死吴军中。军法。父子俱从军。有死事。得与丧归。灌夫不肯随丧归。奋曰。﹁原取吴王若将军头。以报父之仇。﹂於是灌夫被甲持戟。募军中壮士所善原从者数十人。及出壁门。莫敢前。独二人及从奴十数骑驰入吴军。至吴将麾下。所杀伤数十人。不得前。复驰还。走入汉壁。皆亡其奴。独与一骑归。夫身中大创十馀。適有万金良药。故得无死。夫创少瘳。又复请将军曰。﹁吾益知吴壁中曲折。请复往。﹂将军壮义之。恐亡夫。乃言太尉。太尉乃固止之。吴已破。灌夫以此名闻天下。\\
\hspace*{2em}颍阴侯言之上。上以夫为中郎将。数月。坐法去。后家居长安。长安中诸公莫弗称之。孝景时。至代相。孝景崩。今上初即位。以为淮阳天下交。劲兵处。故徙夫为淮阳太守。建元元年。入为太仆。二年。夫与长乐卫尉窦甫饮。轻重不得。夫醉。搏甫。甫。窦太后昆弟也。上恐太后诛夫。徙为燕相。数岁。坐法去官。家居长安。\\
\hspace*{2em}灌夫为人刚直使酒。不好面谀。贵戚诸有势在己之右。不欲加礼。必陵之。诸士在己之左。愈贫贱。尤益敬。与钧。稠人广众。荐宠下辈。士亦以此多之。\\
\hspace*{2em}夫不喜文学。好任侠。已然诺。诸所与交通。无非豪桀大猾。家累数千万。食客日数十百人。陂池田园。宗族宾客为权利。横於颍川。颍川兒乃歌之曰。﹁颍水清。灌氏宁。颍水浊。灌氏族。﹂\\
\hspace*{2em}灌夫家居虽富。然失势。卿相侍中宾客益衰。及魏其侯失势。亦欲倚灌夫引绳批根生平慕之后弃之者。灌夫亦倚魏其而通列侯宗室为名高。两人相为引重。其游如父子然。相得欢甚。无厌。恨相知晚也。\\
\hspace*{2em}灌夫有服。过丞相。丞相从容曰。﹁吾欲与仲孺过魏其侯。会仲孺有服。﹂灌夫曰。﹁将军乃肯幸临况魏其侯。夫安敢以服为解。请语魏其侯帐具。将军旦日蚤临。﹂武安许诺。灌夫具语魏其侯如所谓武安侯。魏其与其夫人益市牛酒。夜洒埽。早帐具至旦。平明。令门下候伺。至日中。丞相不来。魏其谓灌夫曰。﹁丞相岂忘之哉。﹂灌夫不怿。曰。﹁夫以服请。宜往。﹂乃驾。自往迎丞相。丞相特前戏许灌夫。殊无意往。及夫至门。丞相尚卧。於是夫入见。曰。﹁将军昨日幸许过魏其。魏其夫妻治具。自旦至今。未敢尝食。﹂武安鄂谢曰。﹁吾昨日醉。忽忘与仲孺言。﹂乃驾往。又徐行。灌夫愈益怒。及饮酒酣。夫起舞属丞相。丞相不起。夫从坐上语侵之。魏其乃扶灌夫去。谢丞相。丞相卒饮至夜。极欢而去。\\
\hspace*{2em}丞相尝使籍福请魏其城南田。魏其大望曰。﹁老仆虽弃。将军虽贵。宁可以势夺乎。﹂不许。灌夫闻。怒。骂籍福。籍福恶两人有郄。乃谩自好谢丞相曰。﹁魏其老且死。易忍。且待之。﹂已而武安闻魏其。灌夫实怒不予田。亦怒曰。﹁魏其子尝杀人。蚡活之。蚡事魏其无所不可。何爱数顷田。且灌夫何与也。吾不敢复求田。﹂武安由此大怨灌夫。魏其。\\
\hspace*{2em}元光四年春。丞相言灌夫家在颍川。横甚。民苦之。请案。上曰。﹁此丞相事。何请。﹂灌夫亦持丞相阴事。为奸利。受淮南王金与语言。宾客居间。遂止。俱解。\\
\hspace*{2em}夏。丞相取燕王女为夫人。有太后诏。召列侯宗室皆往贺。魏其侯过灌夫。欲与俱。夫谢曰。﹁夫数以酒失得过丞相。丞相今者又与夫有郄。﹂魏其曰。﹁事已解。﹂彊与俱。饮酒酣。武安起为寿。坐皆避席伏。已魏其侯为寿。独故人避席耳。馀半膝席。灌夫不悦。起行酒。至武安。武安膝席曰。﹁不能满觞。﹂夫怒。因嘻笑曰。﹁将军贵人也。属之。﹂时武安不肯。行酒次至临汝侯。临汝侯方与程不识耳语。又不避席。夫无所发怒。乃骂临汝侯曰。﹁生平毁程不识不直一钱。今日长者为寿。乃效女兒呫嗫耳语。﹂武安谓灌夫曰。﹁程李俱东西宫卫尉。今众辱程将军。仲孺独不为李将军地乎。﹂灌夫曰。﹁今日斩头陷匈。何知程李乎。﹂坐乃起更衣。稍稍去。魏其侯去。麾灌夫出。武安遂怒曰。﹁此吾骄灌夫罪。﹂乃令骑留灌夫。灌夫欲出不得。籍福起为谢。案灌夫项令谢。夫愈怒。不肯谢。武安乃麾骑缚夫置传舍。召长史曰。﹁今日召宗室。有诏。﹂劾灌夫骂坐不敬。系居室。遂按其前事。遣吏分曹逐捕诸灌氏支属。皆得弃市罪。魏其侯大媿。为资使宾客请。莫能解。武安吏皆为耳目。诸灌氏皆亡匿。夫系。遂不得告言武安阴事。\\
\hspace*{2em}魏其锐身为救灌夫。夫人谏魏其曰。﹁灌将军得罪丞相。与太后家忤。宁可救邪。﹂魏其侯曰。﹁侯自我得之。自我捐之。无所恨。且终不令灌仲孺独死。婴独生。﹂乃匿其家。窃出上书。立召入。具言灌夫醉饱事。不足诛。上然之。赐魏其食。曰。﹁东朝廷辩之。﹂\\
\hspace*{2em}魏其之东朝。盛推灌夫之善。言其醉饱得过。乃丞相以他事诬罪之。武安又盛毁灌夫所为横恣。罪逆不道。魏其度不可柰何。因言丞相短。武安曰。﹁天下幸而安乐无事。蚡得为肺腑。所好音乐狗马田宅。蚡所爱倡优巧匠之属。不如魏其。灌夫日夜招聚天下豪桀壮士与论议。腹诽而心谤。不仰视天而俯画地。辟倪两宫间。幸天下有变。而欲有大功。臣乃不知魏其等所为。﹂於是上问朝臣。﹁两人孰是。﹂御史大夫韩安国曰。﹁魏其言灌夫父死事。身荷戟驰入不测之吴军。身被数十创。名冠三军。此天下壮士。非有大恶。争杯酒。不足引他过以诛也。魏其言是也。丞相亦言灌夫通奸猾。侵细民。家累巨万。横恣颍川。凌轹宗室。侵犯骨肉。此所谓﹃枝大於本。胫大於股。不折必披﹄。丞相言亦是。唯明主裁之。﹂主爵都尉汲黯是魏其。内史郑当时是魏其。后不敢坚对。馀皆莫敢对。上怒内史曰。﹁公平生数言魏其。武安长短。今日廷论。局趣效辕下驹。吾并斩若属矣。﹂即罢起入。上食太后。太后亦已使人候伺。具以告太后。太后怒。不食。曰。﹁今我在也。而人皆藉吾弟。令我百岁后。皆鱼肉之矣。且帝宁能为石人邪。此特帝在。即录录。设百岁后。是属宁有可信者乎。﹂上谢曰。﹁俱宗室外家。故廷辩之。不然。此一狱吏所决耳。﹂是时郎中令石建为上别言两人事。\\
\hspace*{2em}武安已罢朝。出止车门。召韩御史大夫载。怒曰。﹁与长孺共一老秃翁。何为首鼠两端。﹂韩御史良久谓丞相曰。﹁君何不自喜。夫魏其毁君。君当免冠解印绶归。曰﹃臣以肺腑幸得待罪。固非其任。魏其言皆是﹄。如此。上必多君有让。不废君。魏其必内愧。杜门齰舌自杀。今人毁君。君亦毁人。譬如贾竖女子争言。何其无大体也。﹂武安谢罪曰。﹁争时急。不知出此。﹂\\
\hspace*{2em}於是上使御史簿责魏其所言灌夫。颇不雠。欺谩。劾系都司空。孝景时。魏其常受遗诏。曰﹁事有不便。以便宜论上﹂。及系。灌夫罪至族。事日急。诸公莫敢复明言於上。魏其乃使昆弟子上书言之。幸得复召见。书奏上。而案尚书大行无遗诏。诏书独藏魏其家。家丞封。乃劾魏其矫先帝诏。罪当弃市。五年十月。悉论灌夫及家属。魏其良久乃闻。闻即恚。病痱。不食欲死。或闻上无意杀魏其。魏其复食。治病。议定不死矣。乃有蜚语为恶言闻上。故以十二月晦论弃市渭城。\\
\hspace*{2em}其春。武安侯病。专呼服谢罪。使巫视鬼者视之。见魏其。灌夫共守。欲杀之。竟死。子恬嗣。元朔三年。武安侯坐衣襜褕入宫。不敬。\\
\hspace*{2em}淮南王安谋反觉。治。王前朝。武安侯为太尉。时迎王至霸上。谓王曰。﹁上未有太子。大王最贤。高祖孙。即宫车晏驾。非大王立当谁哉。﹂淮南王大喜。厚遗金财物。上自魏其时不直武安。特为太后故耳。及闻淮南王金事。上曰。﹁使武安侯在者。族矣。﹂\\
\hspace*{2em}太史公曰。魏其。武安皆以外戚重。灌夫用一时决筴而名显。魏其之举以吴楚。武安之贵在日月之际。然魏其诚不知时变。灌夫无术而不逊。两人相翼。乃成祸乱。武安负贵而好权。杯酒责望。陷彼两贤。呜呼哀哉。迁怒及人。命亦不延。众庶不载。竟被恶言。呜呼哀哉。祸所从来矣。\\

% ===== 索隐述赞（司马贞《史记索隐》）· 统一样式 =====
% 题头「索隐述赞曰」五字并入反白框；正文每四句一列，定宽盒 \hfil 居列中；
% \fontsize{24}{41} 列距对齐全书界行（121+41k）；\clearpage 使述赞独占一页。
\clearpage
\shuzanhead{索隐述赞曰}\linebreak
{\shuzanbody\fontsize{24}{41}\selectfont
\hbox to \hsize{\hfil 窦婴。田蚡。势利相雄。咸倚外戚。\hfil}\linebreak
\hbox to \hsize{\hfil 或恃军功。灌夫自喜。引重其中。意气杯酒。\hfil}\linebreak
\hbox to \hsize{\hfil 睥辟睨两宫。事竟不直。冤哉二公。\hfil}%
\par}
\end{document}
